\subsection{模型的优点}
\begin{enumerate}[label=(\arabic*)]
  \item 统一将周期性计划展开为完整重复事件，避免只检查首次时间区间造成漏检。
  \item 用离散时频资源格表达冲突，使 Q2 和 Q3 的可行性能够由独立程序复核。
  \item 对问题二的主要优先级采用条件字典序，并对关键阈值给出上下界或不可行证据，避免把限时可行解直接称为全局最优。
  \item 问题四只增加 C 类间隔调整这一项自由度，并将其通过完整事件递推接入同一资源格模型；三撤销见证、六分支不可行证据和支配覆盖核验使最小撤销层具有可审计的上下界结构。
\end{enumerate}

\subsection{模型的缺点}
\begin{enumerate}[label=(\arabic*)]
  \item 统一时间域右端点 643 是由当前附件展开得到的数据驱动边界，题面没有直接给出该全局终点；虽然进行了 \(H=638,643,648\) 敏感性检验，外推到更宽时间域仍需重新求解。
  \item 问题二末级归一化平移量 \(S_{10}\) 尚未完成全局最小性证明，当前方案的 775 只能作为可行上界。
  \item 问题三的 141 台结果依赖当前 Q2 方案和附件 C 类共同模板，不能直接推广到其他 Q2 背景或其他新增计划结构。
  \item 问题四的 \(R_{Q4}^{\ast}=3\) 只对给定整数动作集合和 \(H=643\) 闭合了最小撤销层；当前 \(R_B=1\)、\(M=142\) 和 \(S_{10}^{(4)}=943\) 尚未完成全部相邻层的不可行证明，不能作为完整字典序最优结论。
\end{enumerate}

\subsection{模型的改进}
后续可在保持资源格精确约束的基础上继续完成 Q2 的 \(S_{10}\) 与 Q4 次级字典序层的阈值证明，并在 \(H=638,648\) 情景下补做最小撤销数的下界核验。若进一步放宽间隔、频率或时间动作，必须重新生成全部候选事件、资源格和跨问接口，而不能在已有工作簿上做局部修补。
